%\subsection{Criterios de selección de herramientas}

%La decisión de conservar o descartar cada herramienta se tomó siguiendo un procedimiento de inclusión/exclusión por herramienta. En primer lugar, se mantuvieron únicamente las soluciones implementadas en el período 2022–2025 cuya ejecución fuera verificable (código, demo o servicio). Se descartaron las herramientas que no pudieron ser ejecutadas o verificadas y los prototipos conceptuales sin implementación funcional. 
%También se eliminaron versiones equivalentes de una misma herramienta, priorizando las alternativas más completas.

%Se conservaron aquellas que abordaran el problema con técnicas de IA/LLMs. En ese marco, el \textbf{criterio principal de inclusión} fue estrictamente funcional: conservar sólo aquellas herramientas que demostraran \emph{generación automática de historias de usuario} a partir de entradas reales del proceso de requisitos, documentos de texto, PDFs y fuentes de audio o imagen.

%Una vez satisfecho el objetivo principal, la continuidad de cada candidata se resolvió aplicando la capa de exclusión/verificación y la evaluación secundaria. En esta instancia, se priorizaron las soluciones que, además de generar historias, permitiesen \emph{mantener la trazabilidad} respecto de los artefactos de origen y exhibieran \emph{capacidad de detectar inconsistencias o ambigüedades} en los documentos de requisitos; ambos aspectos se pusieron a prueba mediante insumos diseñados para observar ese comportamiento y se registraron como evidencia de verificación. En síntesis, una herramienta se conservó si: (i) cumplía el objetivo principal de generación automática desde entradas reales y (ii) superaba la verificación asociada a los aspectos de interés definidos; de lo contrario, se descartó en la etapa de exclusión.

%Para asegurar transparencia y reproducibilidad, la aplicación de este criterio quedó documentada junto con las evidencias de prueba y fechas de ejecución, en línea con los registros metodológicos establecidos para la revisión. 
%(Sugerencia de ref.: añadir aquí una referencia a guías de revisión sistemática en ingeniería de software, p. ej., Kitchenham/Charters.)